\documentclass[11pt]{article}
\usepackage{amsmath,amssymb,booktabs,geometry,graphicx,hyperref,longtable,url}
\usepackage{array,enumitem,mathtools}
\hypersetup{hidelinks}
\geometry{margin=1in}
\setlength{\parskip}{0.55em}
\setlength{\parindent}{0pt}
\setlength{\emergencystretch}{2em}
\setlist[itemize]{topsep=0.25em,itemsep=0.2em}
\setlist[enumerate]{topsep=0.25em,itemsep=0.2em}
\allowdisplaybreaks

\newcommand{\R}{\mathbb{R}}
\newcommand{\payoff}{\Phi}
\newcommand{\stepheading}[1]{\subsubsection*{Step-by-step derivation: #1}}

\newenvironment{readerintuition}
  {\par\medskip\begin{quote}\textbf{Reader intuition.}\ }
  {\end{quote}\medskip}

\newenvironment{commonconfusion}
  {\par\medskip\begin{quote}\textbf{Common confusion.}\ }
  {\end{quote}\medskip}

\newenvironment{implementationcheck}
  {\par\medskip\begin{quote}\textbf{Implementation connection.}\ }
  {\end{quote}\medskip}

\title{Expanded Mathematical Learning Report:\\
American Option Risk Surfaces and Modular Neural Diagnostics}
\author{Free-Boundary-Aware Risk Surfaces for American Options}
\date{\today}

\begin{document}
\maketitle
\begin{abstract}
This report is the learning companion to the concise final report. It explains
the mathematical formulation, finite-difference matrices, CN/PSOR calculation,
boundary and Greek diagnostics, dataset governance, and modular neural
diagnostics for a second-year AI/ML undergraduate. It introduces no new
experiments or numerical claims.
\end{abstract}

\tableofcontents

\section*{How to Use This Learning Version}
\addcontentsline{toc}{section}{How to Use This Learning Version}
Read Stages 1--2 for the mathematical and numerical derivation. Stages 3--4
explain how secondary quantities inherit grid limitations. Stages 5--6 explain
how solver outputs become governed data. Stages 7--11 explain why price,
boundary, and Delta are modeled and audited separately.

\section{Stage 1: Mathematical Foundations}

\subsection{Purpose}
This stage defines the financial problem and the notation used throughout the
project. The objective is not to introduce a new option-pricing theory. It is
to explain why early exercise changes a familiar Black--Scholes equation into
an inequality problem with an unknown boundary.

\subsection{Method}
Let \(S>0\) denote the current asset price, \(K>0\) the strike, \(T\) the
maturity date, \(t\in[0,T]\) calendar time, \(r\) the continuously compounded
risk-free rate, \(q\) the continuous dividend yield, and \(\sigma>0\) the
volatility. The project studies American puts and American calls. Their
immediate-exercise payoffs are
\[
  \payoff_{\mathrm{put}}(S)=\max(K-S,0),
  \qquad
  \payoff_{\mathrm{call}}(S)=\max(S-K,0).
\]
The American value \(V(S,t)\) is the best discounted expected payoff available
when the holder may choose an exercise time \(\theta\) between \(t\) and \(T\):
\[
  V(S,t)
  =
  \sup_{\theta\in[t,T]}
  \mathbb{E}^{\mathbb{Q}}\!
  \left[
    e^{-r(\theta-t)}\payoff(S_\theta)
    \,\middle|\,S_t=S
  \right].
\]
The expectation is under the risk-neutral measure \(\mathbb{Q}\).

\begin{readerintuition}
A European holder has one allowed exercise date. An American holder repeatedly
compares two actions: exercise now for \(\payoff(S)\), or retain the contract.
The option value must therefore be at least the immediate payoff.
\end{readerintuition}

\subsection{Mathematical Formulation}
\subsubsection{Time to maturity}
The numerical solver uses
\[
  \tau=T-t,
  \qquad
  U(S,\tau)=V(S,T-\tau).
\]
Thus \(\tau=0\) is maturity and \(\tau=T\) is the initial valuation date. This
change converts a terminal-value problem in \(t\) into an initial-value problem
in \(\tau\):
\[
  U(S,0)=\payoff(S).
\]

\begin{commonconfusion}
Increasing \(t\) moves forward in calendar time, but increasing \(\tau\) moves
away from maturity. The code begins with the payoff at \(\tau=0\) and advances
through \(\tau_1,\tau_2,\ldots,\tau_N=T\).
\end{commonconfusion}

\subsubsection{Risk-neutral dynamics and the Black--Scholes operator}
Under the model used in the project,
\[
  dS_t=(r-q)S_t\,dt+\sigma S_t\,dW_t.
\]
In a region where exercise is not optimal, standard replication and
risk-neutral valuation give
\[
  U_\tau=\mathcal{L}U,
\]
where
\[
  \mathcal{L}U
  =
  \frac12\sigma^2S^2U_{SS}
  +(r-q)SU_S-rU.
\]
The three terms have separate meanings:
\begin{itemize}
  \item \(\tfrac12\sigma^2S^2U_{SS}\) is the diffusion or curvature term;
  \item \((r-q)SU_S\) is the risk-neutral drift term;
  \item \(-rU\) is the discount term.
\end{itemize}

\subsubsection{From a PDE to a variational inequality}
Immediate exercise is always available, so
\[
  U(S,\tau)\geq\Phi(S).
\]
If \(U(S,\tau)>\payoff(S)\), retaining the option is more valuable and the
continuation PDE applies:
\[
  U_\tau-\mathcal{L}U=0.
\]
If \(U(S,\tau)=\payoff(S)\), the holder exercises. At such a point the PDE
equality need not hold, but delaying exercise cannot produce a locally better
value. The combined problem is
\[
  U(S,\tau)\geq\payoff(S),
  \qquad
  U_\tau-\mathcal{L}U\geq0,
  \qquad
  \bigl(U-\payoff\bigr)
  \bigl(U_\tau-\mathcal{L}U\bigr)=0.
\]

\begin{readerintuition}
The product condition says that at least one factor is zero. In continuation,
\(U-\payoff>0\), so the PDE residual must be zero. In exercise,
\(U-\payoff=0\), while the residual may be positive. This is why American
pricing is an inequality or complementarity problem rather than only a linear
PDE solve.
\end{readerintuition}

\subsubsection{Exercise region, continuation region, and free boundary}
Define
\[
  \mathcal{E}
  =
  \{(S,\tau):U(S,\tau)=\payoff(S)\},
  \qquad
  \mathcal{C}
  =
  \{(S,\tau):U(S,\tau)>\payoff(S)\}.
\]
The interface between \(\mathcal{E}\) and \(\mathcal{C}\) is the free boundary
\(B(\tau)\). It is ``free'' because its location is not prescribed as an input;
it must be discovered together with the option value.

For a put, exercise generally occupies the low-spot side:
\[
  S\lesssim B_{\mathrm{put}}(\tau)
  \quad\Longrightarrow\quad
  \text{exercise}.
\]
For a dividend-paying call, exercise---when it occurs---generally occupies the
high-spot side:
\[
  S\gtrsim B_{\mathrm{call}}(\tau)
  \quad\Longrightarrow\quad
  \text{exercise}.
\]
A no-dividend American call is retained as a control because, under the
standard assumptions, early exercise is not optimal.

\subsubsection{Continuation premium}
The continuation premium is
\[
  C(S,\tau)=U(S,\tau)-\Phi(S).
\]
The obstacle condition immediately gives
\[
  C(S,\tau)\geq0.
\]
Moreover,
\[
  C=0 \quad\text{in the exercise region},
  \qquad
  C>0 \quad\text{in the continuation region}.
\]

\begin{readerintuition}
\(C\) is the extra value of keeping the option alive instead of exercising
immediately. It is therefore both a financial quantity and a useful numerical
indicator of the exercise boundary.
\end{readerintuition}

\begin{commonconfusion}
Nonnegative continuation premium is only the payoff-obstacle requirement. It
does not by itself guarantee monotonicity, convexity, a correct free boundary,
or accurate Delta and Gamma.
\end{commonconfusion}

\subsection{Implementation Details}
The solver stores spot along one array axis and time to maturity along the
other. Reports use normalized quantities such as \(S/K\), \(U/K\), and \(B/K\)
where useful, but the mathematical obstacle remains \(U\geq\payoff\).
Later neural features include
\[
  x=
  \bigl(
    \log(S/K),\tau/T,r,q,\sigma,T,
    \mathbf{1}_{\mathrm{call}}
  \bigr).
\]

\subsection{Validation and Diagnostics}
Before discretization, the notation is checked for the direction of time, the
sign of \(U_\tau-\mathcal{L}U\), put/call payoff orientation, and the units of
\(\Delta=U_S\) and \(\Gamma=U_{SS}\). A sign error can yield stable-looking
arrays while solving the wrong problem.

\subsection{Results}
Stage 1 produces a consistent mathematical contract. It introduces no new
numerical results and makes no model-performance claim.

\subsection{Failure Modes and Limitations}
The formulation assumes the Black--Scholes diffusion with constant parameters.
It does not include jumps, stochastic volatility, transaction costs, or market
calibration. These omissions delimit the project rather than invalidate its
within-model numerical study.

\subsection{Conclusion and Next Step}
The continuous problem is an obstacle-constrained PDE. Stage 2 converts its
spatial operator into a tridiagonal matrix, applies Crank--Nicolson in time, and
solves the resulting linear complementarity problem with PSOR.

\section{Stage 2: From the PDE to CN/PSOR}

\subsection{Purpose}
This stage explains the finite-difference matrix construction in enough detail
to connect every mathematical coefficient with the arrays used by the solver.
The target is a reader who understands derivatives and matrix multiplication
but has not previously implemented an American-option LCP.

\subsection{Method}
\subsubsection{Spot and time grids}
Truncate the unbounded spot domain to \([0,S_{\max}]\). Choose \(M\) equal spot
intervals and \(N\) equal time intervals:
\[
  S_i=i\Delta S,
  \qquad
  i=0,\ldots,M,
  \qquad
  \Delta S=\frac{S_{\max}}{M},
\]
\[
  \tau_n=n\Delta\tau,
  \qquad
  n=0,\ldots,N,
  \qquad
  \Delta\tau=\frac{T}{N}.
\]
Write \(U_i^n\) for the approximation to \(U(S_i,\tau_n)\). At a fixed time,
the interior unknown vector is
\[
  U^n=(U_1^n,U_2^n,\ldots,U_{M-1}^n)^\top.
\]

\begin{commonconfusion}
The index \(i\) moves across spot and the index \(n\) moves across time to
maturity. \(U_i^n\) is one scalar. \(U^n\) is the column vector containing all
interior spot values at time level \(n\).
\end{commonconfusion}

\begin{commonconfusion}
\(U_0^n\) and \(U_M^n\) are not missing unknowns. They are values at the two
domain boundaries \(S=0\) and \(S=S_{\max}\), where an option-specific boundary
condition supplies the value.
\end{commonconfusion}

\subsection{Mathematical Formulation}
\stepheading{central differences}
At an interior node, central differences give
\[
  U_S(S_i,\tau_n)
  \approx
  \frac{U_{i+1}^n-U_{i-1}^n}{2\Delta S},
\]
\[
  U_{SS}(S_i,\tau_n)
  \approx
  \frac{U_{i+1}^n-2U_i^n+U_{i-1}^n}{(\Delta S)^2}.
\]
Substitute them into
\[
  \mathcal{L}U
  =
  \frac12\sigma^2S^2U_{SS}
  +(r-q)SU_S-rU.
\]
At \(S_i\),
\begin{align*}
  (L_hU)_i
  &=
  \frac12\sigma^2S_i^2
  \frac{U_{i+1}-2U_i+U_{i-1}}{(\Delta S)^2}\\
  &\quad
  +(r-q)S_i
  \frac{U_{i+1}-U_{i-1}}{2\Delta S}
  -rU_i.
\end{align*}
Collect the coefficients of the three neighboring values:
\[
  (L_hU)_i=a_iU_{i-1}+b_iU_i+c_iU_{i+1},
\]
where
\[
  a_i=
  \frac{\tfrac12\sigma^2S_i^2}{(\Delta S)^2}
  -\frac{(r-q)S_i}{2\Delta S},
\]
\[
  b_i=
  -\frac{\sigma^2S_i^2}{(\Delta S)^2}-r,
\]
\[
  c_i=
  \frac{\tfrac12\sigma^2S_i^2}{(\Delta S)^2}
  +\frac{(r-q)S_i}{2\Delta S}.
\]

\subsubsection{Where the coefficients come from}
\begin{center}
\begin{tabular}{>{\raggedright\arraybackslash}p{0.18\textwidth}
                >{\raggedright\arraybackslash}p{0.23\textwidth}
                >{\raggedright\arraybackslash}p{0.23\textwidth}
                >{\raggedright\arraybackslash}p{0.23\textwidth}}
\toprule
Source & Contribution to \(a_i\) & Contribution to \(b_i\) & Contribution to \(c_i\)\\
\midrule
Diffusion
& \(+\tfrac12\sigma^2S_i^2/(\Delta S)^2\)
& \(-\sigma^2S_i^2/(\Delta S)^2\)
& \(+\tfrac12\sigma^2S_i^2/(\Delta S)^2\)\\
Drift
& \(-(r-q)S_i/(2\Delta S)\)
& \(0\)
& \(+(r-q)S_i/(2\Delta S)\)\\
Discount
& \(0\)
& \(-r\)
& \(0\)\\
\bottomrule
\end{tabular}
\end{center}

\begin{readerintuition}
Diffusion spreads weight symmetrically to the left and right. Drift shifts the
two off-diagonal weights in opposite directions. Discounting affects only the
current node and therefore appears only in the diagonal coefficient.
\end{readerintuition}

\stepheading{a small matrix with \(M=4\)}
With nodes \(S_0,S_1,S_2,S_3,S_4\), there are three interior unknowns:
\[
  U^n=(U_1^n,U_2^n,U_3^n)^\top.
\]
The interior matrix is
\[
  L_h=
  \begin{bmatrix}
    b_1 & c_1 & 0\\
    a_2 & b_2 & c_2\\
    0 & a_3 & b_3
  \end{bmatrix}.
\]
The full stencil equations are
\begin{align*}
  (L_hU)_1&=a_1U_0+b_1U_1+c_1U_2,\\
  (L_hU)_2&=a_2U_1+b_2U_2+c_2U_3,\\
  (L_hU)_3&=a_3U_2+b_3U_3+c_3U_4.
\end{align*}
Matrix multiplication by the displayed \(3\times3\) matrix contains only the
interior terms. The omitted \(a_1U_0\) and \(c_3U_4\) are known boundary
contributions.

\begin{commonconfusion}
The first matrix row has no column for \(U_0\), and the last row has no column
for \(U_4\), because the matrix maps the three-component interior vector to
another three-component vector. Only the first stencil touches the lower
boundary; only the last stencil touches the upper boundary.
\end{commonconfusion}

\stepheading{Crank--Nicolson time stepping}
Crank--Nicolson averages the spatial operator at the old and new time levels:
\[
  \frac{U^{n+1}-U^n}{\Delta\tau}
  =
  \frac12L_hU^{n+1}
  +\frac12L_hU^n.
\]
Multiply by \(\Delta\tau\):
\[
  U^{n+1}-U^n
  =
  \frac{\Delta\tau}{2}L_hU^{n+1}
  +\frac{\Delta\tau}{2}L_hU^n.
\]
Move the new-time operator to the left:
\[
  U^{n+1}
  -\frac{\Delta\tau}{2}L_hU^{n+1}
  =
  U^n
  +\frac{\Delta\tau}{2}L_hU^n.
\]
Define
\[
  A=I-\frac{\Delta\tau}{2}L_h,
  \qquad
  B=I+\frac{\Delta\tau}{2}L_h.
\]
Then
\[
  AU^{n+1}=BU^n+g^{n,n+1},
\]
where \(g^{n,n+1}\) collects known boundary values.

\begin{commonconfusion}
\(A\) multiplies \(U^{n+1}\) because its terms were gathered on the left.
\(B\) multiplies \(U^n\) because its terms remain on the right. The signs are a
consequence of rearrangement, not an arbitrary naming convention.
\end{commonconfusion}

\subsubsection{What one row means}
Ignoring the boundary adjustment momentarily, an interior row \(i\) reads
\begin{align*}
 &-\frac{\Delta\tau}{2}a_iU_{i-1}^{n+1}
 +\left(1-\frac{\Delta\tau}{2}b_i\right)U_i^{n+1}
 -\frac{\Delta\tau}{2}c_iU_{i+1}^{n+1}\\
 &\qquad=
 \frac{\Delta\tau}{2}a_iU_{i-1}^{n}
 +\left(1+\frac{\Delta\tau}{2}b_i\right)U_i^{n}
 +\frac{\Delta\tau}{2}c_iU_{i+1}^{n}.
\end{align*}
This is the scalar equation represented by row \(i\) of \(A\) and \(B\).

\stepheading{why boundary values appear on the right-hand side}
For the first interior row, the lower-boundary terms at the two time levels are
\[
  a_1U_0^{n+1},
  \qquad
  a_1U_0^n.
\]
Moving the new-time boundary term to the right gives the combined contribution
\[
  \frac{\Delta\tau}{2}a_1
  \left(U_0^n+U_0^{n+1}\right).
\]
Similarly, the last row receives
\[
  c_{M-1}U_M^{n+1},
  \qquad
  c_{M-1}U_M^n,
\]
which combine as
\[
  \frac{\Delta\tau}{2}c_{M-1}
  \left(U_M^n+U_M^{n+1}\right).
\]
Hence
\[
  g^{n,n+1}
  =
  \frac{\Delta\tau}{2}
  \begin{bmatrix}
    a_1(U_0^n+U_0^{n+1})\\
    0\\
    \vdots\\
    0\\
    c_{M-1}(U_M^n+U_M^{n+1})
  \end{bmatrix}.
\]

\begin{implementationcheck}
This matches \path{src/solvers/cn_psor.py}: the first right-hand-side entry is
increased by
\(\tfrac12\Delta\tau\,a_1(U_0^n+U_0^{n+1})\), and the last by
\(\tfrac12\Delta\tau\,c_{M-1}(U_M^n+U_M^{n+1})\).
\end{implementationcheck}

\subsubsection{Boundary conditions used in the project}
For an American put,
\[
  U(0,\tau)=K,
  \qquad
  U(S_{\max},\tau)=0.
\]
At zero spot, immediate exercise returns \(K\). At a sufficiently large spot,
the put value approaches zero.

For a dividend-paying American call,
\[
  U(0,\tau)=0,
\]
\[
  U(S_{\max},\tau)
  =
  \max\!\left(
    S_{\max}-K,\,
    S_{\max}e^{-q\tau}-Ke^{-r\tau},\,
    0
  \right).
\]
The upper boundary retains both immediate-exercise value and the discounted
European far-field expression.

For European validation with projection disabled, the put uses
\[
  U(0,\tau)=Ke^{-r\tau},
  \qquad
  U(S_{\max},\tau)=0,
\]
and the call uses
\[
  U(0,\tau)=0,
  \qquad
  U(S_{\max},\tau)=S_{\max}e^{-q\tau}-Ke^{-r\tau}.
\]

\stepheading{the American linear complementarity problem}
Let
\[
  d^n=BU^n+g^{n,n+1}.
\]
A European step solves the equality \(AU^{n+1}=d^n\). For an American option,
that equality can produce a value below the payoff. The discrete obstacle
problem is instead
\[
  AU^{n+1}\geq d^n,
  \qquad
  U^{n+1}\geq\payoff,
\]
\[
  \left(U^{n+1}-\payoff\right)^\top
  \left(AU^{n+1}-d^n\right)=0.
\]
The first inequality requires a nonnegative discrete PDE residual. The second
enforces the payoff obstacle. The final equality imposes complementarity:
each node is either continuing with zero PDE residual or exercising with zero
premium.

\begin{commonconfusion}
An American option is not priced by solving the European system and then
reporting whichever nodes happen to lie above the payoff. The obstacle is part
of the system itself, so the solution method must enforce it during the solve.
\end{commonconfusion}

\stepheading{projected successive over-relaxation}
At one time step, write \(Ax=d\) for the interior system and let
\(\payoff_i\) be the obstacle at node \(i\). Starting from an old iterate
\(x^{(k)}\), a Gauss--Seidel candidate is
\[
  y_i^{(k+1)}
  =
  \frac{
    d_i
    -\sum_{j<i}A_{ij}x_j^{(k+1)}
    -\sum_{j>i}A_{ij}x_j^{(k)}
  }{A_{ii}}.
\]
The first sum uses already updated entries; the second uses entries from the
previous sweep. Relaxation gives
\[
  x_i^{\mathrm{relaxed}}
  =
  x_i^{(k)}
  +\omega\left(y_i^{(k+1)}-x_i^{(k)}\right).
\]
Projection then enforces the obstacle:
\[
  x_i^{\mathrm{new}}=\max\!\left(\payoff_i,
  x_i^{\mathrm{relaxed}}\right).
\]
The parameter \(\omega\in(0,2)\) controls relaxation. Values above one may
accelerate convergence for suitable systems, but excessive relaxation can
destabilize the iteration.

\subsubsection{Stopping and complementarity diagnostics}
The implemented iteration stops when
\[
  \lVert x^{(k+1)}-x^{(k)}\rVert_\infty<\text{tol},
\]
or reports failure after the maximum number of sweeps. A small update is
necessary but is not the full financial diagnostic. After solving, define
\[
  z=Ax-d.
\]
The LCP checks are
\[
  x-\payoff\geq0,
  \qquad
  z\geq0,
  \qquad
  (x-\payoff)^\top z\approx0.
\]
These separately test obstacle feasibility, residual feasibility, and
complementarity.

\begin{readerintuition}
PSOR alternates two actions: move toward the linear-system solution, then clip
the result back to the financially admissible side of the payoff. Convergence
is credible only when both the numerical update and the LCP residual checks are
small.
\end{readerintuition}

\subsection{Implementation Details}
The tridiagonal operator is exposed by
\path{src/solvers/cn_psor.py} for sign and matrix tests. Boundary formulas live
in \path{src/solvers/boundary_conditions.py}. The code uses the payoff as the
initial condition at \(\tau=0\), forms \(A\) and \(B\) once per solve, adjusts
the first and last right-hand-side entries at every time step, and applies PSOR
only when American projection is enabled.

\subsection{Validation and Diagnostics}
The existing solver audit tests verify the operator on a linear function, the
signs in \(A\) and \(B\), put and dividend-call boundary values, and PSOR
projection/complementarity. European closed-form and no-dividend call controls
provide independent checks beyond internal matrix consistency.

\subsection{Results}
The verified smoke-grid maximum error is \(2.782\times10^{-4}\) for both the
European put validation and the no-dividend American call control. The American
put obstacle violation is reported as zero, with complementarity error
\(4.009\times10^{-12}\). These are engineering-scale validation results, not a
claim of universal or publication-grade accuracy.

\subsection{Failure Modes and Limitations}
Uniform grids resolve the payoff kink and free boundary only approximately.
Finite-domain truncation can contaminate the interior if \(S_{\max}\) is too
small. Crank--Nicolson begins from a nonsmooth payoff, and Rannacher smoothing
is not active in the current solver. These limitations are particularly
important for Gamma near maturity.

\subsection{Conclusion and Next Step}
The PDE has become a boundary-adjusted LCP at each time step. Later stages
extract the exercise boundary and Greeks, study resolution and runtime, and use
only diagnostic-passing solver outputs for machine-learning labels.

\section{Stage 3: Boundary Extraction and Greeks}

\subsection{Purpose and Method}
Once the CN/PSOR surface is available, Stage 3 derives secondary quantities
without pretending that they are as numerically stable as price. For every time
level, compute
\[
  C_i^n=U_i^n-\Phi(S_i).
\]
Nodes with \(C_i^n\) at or below a small tolerance are classified as exercise
nodes; larger premiums indicate continuation.

\subsection{Boundary orientation and status}
For a put, the exercise set begins at low spot. The algorithm finds the first
continuation node and interpolates between it and the preceding exercise node.
For a dividend-paying call, exercise is on the high-spot side. The algorithm
finds the last continuation node and interpolates toward the next exercise
node. These directions are deliberately different.

The output records one of three statuses:
\begin{itemize}
  \item \texttt{boundary\_found}: both exercise and continuation nodes exist;
  \item \texttt{no\_exercise}: no exercise node appears on the truncated grid;
  \item \texttt{all\_exercise}: every grid node is classified as exercise.
\end{itemize}
If no boundary exists, the code records a missing value rather than fabricating
a threshold.

\begin{commonconfusion}
A no-dividend American call control should generally have no early-exercise
boundary. Replacing its missing boundary with zero or \(S_{\max}\) would turn a
valid theorem-based control into a false supervised label.
\end{commonconfusion}

\subsection{Delta, Gamma, and strict masks}
For interior nodes, the teaching approximations are
\[
  \Delta_i^n
  \approx
  \frac{U_{i+1}^n-U_{i-1}^n}{2\Delta S},
\]
\[
  \Gamma_i^n
  \approx
  \frac{U_{i+1}^n-2U_i^n+U_{i-1}^n}{(\Delta S)^2}.
\]
Delta is a first derivative; Gamma is a second derivative. Differentiation
amplifies grid and interpolation error, and the second derivative amplifies it
again. The payoff kink at \(S=K\) and the exercise transition at \(S=B(\tau)\)
are therefore delicate regions.

The strict Greek mask excludes nodes near zero spot, the payoff kink, and a
finite band around the free boundary. The mask does not make Gamma exact. It
defines a less singular region in which comparisons are more meaningful.

\begin{readerintuition}
Price can look smooth and accurate while Gamma oscillates, because many nearby
price errors cancel in the level but become magnified by repeated
differencing.
\end{readerintuition}

\section{Stage 4: Grid, Domain, and Runtime}

\subsection{Why resolution matters}
A finite-difference answer depends on \(\Delta S\), \(\Delta\tau\), and
\(S_{\max}\). Refining the grid reduces discretization error but increases
matrix size and the number of time steps. Boundary locations are especially
grid-limited because the exercise transition is first detected between
neighboring spot nodes.

The grid-240 solution is used as an internal numerical reference. It is not an
analytic truth surface. If \(Q_h\) is a quantity computed on grid \(h\), the
reported comparison has the form
\[
  \mathrm{RMSE}(Q_h,Q_{\mathrm{ref}})
  =
  \left(
    \frac{1}{|\Omega|}
    \sum_{(i,n)\in\Omega}
    \left(Q_{h,i}^n-Q_{\mathrm{ref},i}^n\right)^2
  \right)^{1/2},
\]
after interpolation where required. For Greeks, \(\Omega\) is restricted by
the strict mask.

\subsection{Domain truncation}
The true Black--Scholes spot domain is unbounded, whereas the computation stops
at \(S_{\max}\). If \(S_{\max}\) is too small, the imposed upper boundary can
affect economically relevant interior nodes. Comparing several multiples of
\(K\) tests this sensitivity but does not prove domain independence for every
future parameter regime.

\subsection{Convergence and runtime interpretation}
Price generally converges more smoothly than the free boundary, Delta, or
Gamma. Gamma converges most slowly because it is a second derivative and is
sensitive to nonsmooth initial data. Runtime rises with the number of time
steps, the number of spatial unknowns, and the PSOR sweeps required at each
step. The project reports runtime together with iteration counts so that speed
is not separated from solver convergence.

\begin{commonconfusion}
A small error relative to grid 240 means agreement with a finer internal
calculation. It does not mean agreement with an exact American-option formula,
because no such general closed form is being used as the reference.
\end{commonconfusion}

\section{Stages 5--6: Risk Surfaces and Dataset Governance}

\subsection{What a risk surface represents}
A risk surface evaluates price or a diagnostic over many spot, time, and model
parameter combinations rather than at one contract point. This reveals how
early exercise and sensitivities vary across regimes.

The application sweep separates:
\begin{itemize}
  \item American puts, whose exercise region is on the low-spot side;
  \item dividend-paying American calls, which may exercise on the high-spot
        side;
  \item no-dividend American call controls, which test the no-early-exercise
        theorem rather than provide dividend-call boundary evidence.
\end{itemize}
The verified application set contains 288 regimes: 144 puts, 108
dividend-paying calls, and 36 no-dividend controls. Thirty-two regimes carry
diagnostic failures and remain visible in the audit.

\subsection{Audit data versus clean labels}
Audit retention and model-label eligibility serve different purposes. A failed
or delicate regime remains scientifically important because it shows where the
solver diagnostics or domain assumptions need attention. It is excluded from
the clean machine-learning labels because training on a questionable target
would transfer numerical failure into the surrogate.

After controls and flagged regimes are excluded, the clean dataset contains
223 regimes and 3,264,943 samples. The 65 excluded regimes are recorded rather
than deleted from the project history.

\subsection{Why the split is by regime}
All spot-time nodes generated from one parameter tuple share the same option
surface. A pointwise random split could place neighboring points from the same
surface in both training and test sets. The model could then appear to
generalize while mostly interpolating within a regime it has already seen.

A regime-level split assigns the entire surface to one partition. The verified
split contains 156 training regimes, 33 validation regimes, and 34 test
regimes.

\begin{readerintuition}
Treat one parameter regime like one patient in a medical dataset: many
measurements from the same underlying case should not be divided between
training and testing.
\end{readerintuition}

\section{Stages 7--8: Direct Price and Positive Premium}

\subsection{Why low RMSE is insufficient}
A direct-price regressor learns
\[
  \widehat U_\theta=f_\theta(x).
\]
Standard squared-error training rewards average price accuracy but does not
encode \(\widehat U\geq\Phi\). A model can therefore achieve low RMSE while
predicting below the immediate payoff at some nodes.

A soft obstacle penalty may add
\[
  \lambda
  \left[\max\!\left(\Phi-\widehat U,0\right)\right]^2
\]
to the loss. This discourages violations but remains a trade-off with the data
loss. Finite \(\lambda\), finite optimization, and unseen inputs mean the
penalty is not a hard guarantee.

\subsection{Positive continuation-premium reconstruction}
Instead, predict a nonnegative premium:
\[
  \widehat C_\theta=\operatorname{softplus}(g_\theta(x))\geq0,
\]
\[
  \widehat U=\Phi+\widehat C.
\]
Therefore
\[
  \widehat U=\Phi+\widehat C,\quad \widehat C\geq0
  \Rightarrow
  \widehat U\geq\Phi.
\]
This removes one specific failure mode: payoff-obstacle violation.

\begin{commonconfusion}
Obstacle consistency is not full financial correctness. The construction does
not enforce monotonicity, convexity, boundary accuracy, Delta accuracy, or
Gamma accuracy. It also does not make the neural model a substitute for the
CN/PSOR benchmark.
\end{commonconfusion}

The selected positive-premium MLP has validation RMSE \(0.0200\), test RMSE
\(0.0218\), and zero obstacle violation on evaluated samples. The linear
premium sanity baseline has lower aggregate RMSE in this run, so the neural
model is not claimed to dominate all baselines.

\section{Stages 9--10: Boundary and Delta Diagnostics}

\subsection{Why boundary prediction is difficult}
The free boundary is an implicit transition derived from an entire option
surface. It changes orientation by option type, may be absent, and is limited
by the finite-difference grid. Boundary labels are also repeated over many
spot rows for the same regime-time pair, so a large point count does not imply
an equally large number of independent boundaries.

The boundary model is therefore a separate supervised diagnostic head:
\[
  \widehat B/K=h_\theta(\tau/T,r,q,\sigma,T,\text{option type}).
\]
It does not solve the obstacle PDE or reproduce PSOR iterations. The pooled
boundary MLP test RMSE is \(0.3513\) and its test MAE is \(0.2087\), making it
the weakest of the three neural components.

\subsection{Why price autograd Delta requires validation}
Let
\[
  y=\frac{U}{K},
  \qquad
  z=\log(S/K).
\]
Because \(U=Ky\) and
\[
  \frac{\partial z}{\partial S}=\frac{1}{S},
\]
the chain rule gives
\[
  \Delta
  =
  \frac{\partial U}{\partial S}
  =
  K\frac{\partial y}{\partial z}\frac{1}{S}
  =
  \frac{\partial (U/K)}{\partial \log(S/K)}
  \frac{1}{S/K}.
\]
A differentiable price network supplies this derivative computationally, but
price-loss accuracy does not directly train the derivative. Small level errors
can still create material slope errors.

The price-model autograd Delta test RMSE is \(0.0977\), which motivates a
separate supervised head. Its output map imposes
\[
  \Delta_{\mathrm{put}}\in[-1,0],
  \qquad
  \Delta_{\mathrm{call}}\in[0,1].
\]
The bounded Delta MLP has strict-mask test RMSE \(0.0332\), test MAE \(0.0140\),
and zero bound violation on evaluated samples.

\begin{commonconfusion}
Bounding Delta limits its range; it does not prove that the predicted slope is
correct. It also gives no second-derivative control and therefore does not
support a Gamma surrogate.
\end{commonconfusion}

\section{Stage 11: Integrated Scenario Screening}

\subsection{Why the workflow has three heads}
The integrated workflow keeps separate mathematical responsibilities:
\begin{enumerate}
  \item the positive-premium head estimates price while enforcing the payoff
        obstacle;
  \item the boundary head estimates a PSOR-derived transition as a diagnostic;
  \item the bounded Delta head estimates a first derivative within
        option-specific ranges.
\end{enumerate}
Combining their outputs in one interface does not turn them into one universal
model. Each head has separate labels, losses, metrics, and limitations.

\subsection{Scenario screening, not replacement}
Scenario screening means generating rapid preliminary estimates for cases
inside the tested regime envelope. Fresh CN/PSOR remains required for new
regimes, final label generation, and audit-quality decisions. The integrated
experiment evaluates six deterministic held-out regimes: three puts and three
dividend-paying calls.

Across those cases, the verified mean price RMSE is \(0.0092\), mean boundary
MAE is \(0.1539\), and mean strict-mask Delta RMSE is \(0.0300\). Obstacle and
Delta-bound violations are zero on the audited cases. Six cases, all from the
designed test envelope, are not enough to establish universal generalization.

\begin{readerintuition}
The workflow is a fast screening layer followed by a numerical audit, much as
a preliminary model can prioritize scenarios before a higher-fidelity
simulation is run.
\end{readerintuition}

\subsection{Overall Failure Modes and Limitations}
The expanded workflow retains the original conservative boundaries:
\begin{itemize}
  \item grid 240 is an internal reference, not analytic truth;
  \item Rannacher smoothing is not active;
  \item boundary estimation remains weaker than price and Delta estimation;
  \item strict masks exclude the most delicate Greek regions;
  \item the evidence does not support a Gamma surrogate;
  \item the neural workflow does not replace CN/PSOR;
  \item six integrated audit regimes do not establish broad generalization.
\end{itemize}

\subsection{Conclusion and Next Step}
The project demonstrates a benchmark-first learning workflow: formulate the
obstacle problem, validate CN/PSOR, attach diagnostics, govern labels, then
evaluate modular neural approximations under explicit claim limits. Stronger
claims would require broader independent regimes, higher-resolution numerical
references, and additional derivative validation.


\end{document}
